(1) 代入平面波，位涡因子 \((\nabla^2-L_d^{-2})\psi=-(k_x^2+k_y^2+L_d^{-2})\psi\)。于是
\[
\boxed{\omega=-\frac{\beta k_x}{K^2}},\quad K^2=k_x^2+k_y^2+L_d^{-2}.
\]
对 \(k_x>0\)，东西向相速度 \(\omega/k_x=-\beta/K^2<0\)，即相位整体西传。有限形变半径把波数平方中的零点抬高，防止极长波的奇异传播速度；这一低频Rossby恢复力来自 \(\beta\)，而非均匀 \(f\) 下的快惯性重力波。

(2) 直接对色散微分得
\[
c_{gx}=-\beta\frac{K^2-2k_x^2}{K^4},\qquad
c_{gy}=2\beta\frac{k_xk_y}{K^4}.
\]
对于 \(k_y=0\)，\(c_{gx}=\beta(k_x^2-L_d^{-2})/(k_x^2+L_d^{-2})^2\)。因此当 \(k_xL_d>1\) 能量包络可东传，而相速度始终西向；\(k_xL_d<1\) 则两者都西传。能量传输、位涡波的相位面移动与实际物质平流不可混为一个速度。

(3) 令 \(Q=\nabla^2\psi-L_d^{-2}\psi\)，方程为 \(Q_t=-\beta\psi_x\)。定义
\[
\mathcal E=\frac12\int\bigl(|\nabla\psi|^2+L_d^{-2}\psi^2\bigr)\,dA=-\frac12\int\psi Q\,dA .
\]
周期积分分部并利用自伴的 \(Q\) 算符可得
\(\dot{\mathcal E}=-\int\psi Q_t\,dA=\beta\int\psi\psi_x\,dA=0\)。单一非零Fourier实模平均 \(\langle\mathcal E\rangle/A_{\rm box}=|A|^2K^2/4>0\)，这是对色散和正定性相互独立的核验。

(4) \(L_d\to\infty\) 时 \(\omega=-\beta k_x/(k_x^2+k_y^2)\)（非零波数）；\(\beta\to0\) 时 \(\omega\to0\)，说明该恢复力随纬度梯度消失。长波 \(kL_d\ll1\) 有 \(\omega=-\beta L_d^2k_x[1+O(k^2L_d^2)]\)，表现近无色散的西向传播。常数 \(f\) 的惯性重力波满足不同的 \(\omega^2=f^2+c^2k^2\)，故不能从其平方色散套用本题的群速度。
